%  LaTeX support: latex@mdpi.com 
%  For support, please attach all files needed for compiling as well as the log file, and specify your operating system, LaTeX version, and LaTeX editor.

%=================================================================
\documentclass[land,article,submit,pdftex,moreauthors]{Definitions/mdpi} 

%=================================================================
%----------
% journal
%----------

% MDPI internal commands
\firstpage{1} 
\makeatletter 
\setcounter{page}{\@firstpage} 
\makeatother
\pubvolume{1}
\issuenum{1}
\articlenumber{0}
\pubyear{2023}
\copyrightyear{2023}
%\externaleditor{Academic Editor: Firstname Lastname}
\datereceived{31 July 2023} 
\daterevised{}
\dateaccepted{} 
\datepublished{} 
\hreflink{https://doi.org/} % If needed use \linebreak
%\doinum{}

%=================================================================
% Add packages and commands here. The following packages are loaded in our class file: fontenc, inputenc, calc, indentfirst, fancyhdr, graphicx, epstopdf, lastpage, ifthen, lineno, float, amsmath, setspace, enumitem, mathpazo, booktabs, titlesec, etoolbox, tabto, xcolor, soul, multirow, microtype, tikz, totcount, changepage, attrib, upgreek, cleveref, amsthm, hyphenat, natbib, hyperref, footmisc, url, geometry, newfloat, caption

\graphicspath{{figures/}} % path to folder with figures
\usepackage{subcaption}
\usepackage{listings}
\usepackage{xcolor}
\usepackage{gensymb} % for \textdegree
\usepackage{tabularx}

%=================================================================
% Full title of the paper (Capitalized)
\Title{Mapping Dynamics of the Land Cover Types in Senegal, West Africa Using Image Synthesis and Analysis by GRASS GIS}

% MDPI internal command: Title for citation in the left column
\TitleCitation{Mapping Dynamics of the Land Cover Types in Senegal, West Africa Using Image Synthesis and Analysis by GRASS GIS}

% Author Orchid ID: enter ID or remove command
\newcommand{\orcidauthorA}{0000-0002-5759-1089} % Polina Add \orcidA{} behind the author's name

% Authors, for the paper (add full first names)
\Author{Polina Lemenkova $^{1}$*\orcidA{}}

%\longauthorlist{yes}

% MDPI internal command: Authors, for metadata in PDF
\AuthorNames{Polina Lemenkova}

% MDPI internal command: Authors, for citation in the left column
\AuthorCitation{Lemenkova, P.}
% If this is a Chicago style journal: Lastname, Firstname, Firstname Lastname, and Firstname Lastname.

% Affiliations / Addresses (Add [1] after \address if there is only one affiliation.)
\address{%
$^{1}$ \quad Laboratory of Image Synthesis and Analysis (LISA), École polytechnique de Bruxelles (Brussels Faculty of Engineering), Université Libre de Bruxelles (ULB). Building L, Campus du Solbosch, ULB -- LISA CP165/57, Avenue Franklin D. Roosevelt 50, Brussels 1050, Belgium.
}

% Contact information of the corresponding author
\corres{Correspondence: polina.lemenkova@ulb.be; Tel.: +32 471 86 04 59 (P.L.)}

% Current address and/or shared authorship
%\firstnote{Current address: Affiliation 3.} 

% Abstract (Do not insert blank lines, i.e. \\) 
\abstract{This paper addresses the problem of mapping land cover types in Senegal and recognition of vegetation systems in the Saloum Delta River on the satellite images. In monitoring of landscape dynamics, the goal of the remote sensing approach is to find all the identical landscape patterns on a time series of images and to analyse changes using recognition of landscape patches. Multi-seasonal landscape dynamics in the region of Saloum River was analysed on the basis of Landsat 8-9 OLI/TIRS datasets on 2014 to 2023 and GRASS GIS scripts using image analysis techniques. Contrary to the existing GIS, programming approach to satellite image processing has a high automation in data handling and is very fast and robust to classify and to compare the satellite images. This is achieved by a combination of the embedded GRASS GIS modules used separately for computing the vegetation indices and plotting the classification maps. The land cover types were identified based on the analysis of attributed of the spectral reflectance of the multispectral images. The detected land cover changes indicate climate change factors and the effects of environmental disasters in semi-arid region of West Africa: variability of precipitation, increase in temperature, land erosions and wildfire. These factors control the distribution of croplands, natural vegetation systems and coastal mangrove colonies in western Senegal. The results based on the observed satellite images indicated the decrease in savannahs, increase in croplands and agricultural lands, decline in forests and changes in coastal wetlands including mangroves with high biodiversity. Mapping land cover types, discriminating mangroves and monitoring savannah ecosystems is essential for proper land management of Senegal. This paper contributed to the natural resource management and environmental monitoring of Senegal through applied advanced cartographic methods to remote sensing Earth observation data on Senegal, West Africa.
}

% Keywords
\keyword{Remote sensing; Cartography; Vegetation; West Africa; Satellite Image; Landsat; Sahel; Climate Change; keyword 9; keyword 10} 

% The fields PACS, MSC, and JEL may be left empty or commented out if not applicable
%\PACS{J0101}
%\MSC{}
%\JEL{}

\begin{document}

%----------- section ----------->
\section{Introduction}

Text. 

\begin{figure}[H]
	\centering
		\includegraphics[width=15.0 cm]{Fig_01.jpg}
\caption{Study area with shown segments of the Landsat images on a topographic map of Senegal. Software: GMT. Map source: author.}\label{fig01}
\end{figure}

%----------- section ----------->
\section{Materials and Methods}

%----------- subsection ----------->
\subsection{Subsection}

%----------- subsection ----------->
\subsection{Subsection}

%----------- subsection ----------->
\subsection{Subsection}

%----------- section ----------->
\section{Results}

%----------- section ----------->
\section{Discussion}

%----------- section ----------->
\section{Conclusions}

%----------- section ----------->
\authorcontributions{Supervision, conceptualization, methodology, software, resources, funding acquisition, and project administration, writing---original draft preparation, methodology, software, data curation, visualization, formal analysis, validation, writing---review and editing, and investigation, P.L.}

\funding{The publication was funded by the Editorial Office of Land, Multidisciplinary Digital Publishing Institute (MDPI), by providing 100\% discount for the APC of this manuscript.}

\institutionalreview{Not applicable.}

\informedconsent{Not applicable.}

\dataavailability{Not applicable} 

\acknowledgments{The authors thank the reviewers for reading and review of this manuscript.}

\conflictsofinterest{The authors declare no conflict of interest.} 

\abbreviations{Abbreviations}{
The following abbreviations are used in this manuscript:\\

\noindent 
\begin{tabular}{@{}ll}
DCW & Digital Chart of the World \\
DEM & Digital Elevation Model \\ 
GEBCO & General Bathymetric Chart of the Oceans \\ 
GMT & Generic Mapping Tools \\
GRASS & Geographic Resources Analysis Support System\\
GIS & Geographic Information System \\
Landsat OLI/TIRS & Landsat Operational Land Imager and Thermal Infrared Sensor \\
TIFF & Tag Image File Format\\
USGS & United States Geological Survey\\
\end{tabular}
}

%%%%%%%%%%%%%%%%%%%%%%%%%%%%%%%%%%%%%%%%%%
%% Optional
\appendixtitles{no} % Leave argument "no" if all appendix headings stay EMPTY (then no dot is printed after "Appendix A"). If the appendix sections contain a heading then change the argument to "yes".
\appendixstart
\appendix
\section[\appendixname~\thesection]{}
\subsection[\appendixname~\thesubsection]{}
The appendix 

\section[\appendixname~\thesection]{}
All appendix sections must be cited in the main text. In the appendices, Figures, Tables, etc. should be labeled, starting with ``A''---e.g., Figure A1, Figure A2, etc.

%%%%%%%%%%%%%%%%%%%%%%%%%%%%%%%%%%%%%%%%%%
\begin{adjustwidth}{-\extralength}{0cm}
%\printendnotes[custom] % Un-comment to print a list of endnotes

\reftitle{References}

\bibliography{Senegal.bib}
\nocite{*}

%%%%%%%%%%%%%%%%%%%%%%%%%%%%%%%%%%%%%%%%%%
\end{adjustwidth}
\end{document}
